% @@STATS_BEGIN@@ — auto-generated by habor-analyze, do not edit manually
% Auto-generated by habor-analyze — do not edit manually.
% Re-generate: uv run habor-analyze studies

% ── Appendix: BenchPress general ──
\providecommand{\AppNumSystems}{16}
\providecommand{\AppNumBenchmarks}{54}
\providecommand{\AppNumModelsOnly}{8}

% ── Finding 1: SVD spectrum ──
\providecommand{\AppSVDModelOnlyPCOne}{71.8}
\providecommand{\AppSVDModelOnlyPCOneTwo}{80.4}
\providecommand{\AppSVDFullPCOne}{65.2}
\providecommand{\AppSVDFullPCTwo}{10.5}
\providecommand{\AppSVDFullCumTwo}{75.7}

% ── Finding 2: PC2 agent separation ──
\providecommand{\AppNumAgentSystems}{8}
\providecommand{\AppNumBaseSystems}{8}
\providecommand{\AppPCTwoAgentMean}{+1.83}
\providecommand{\AppPCTwoBaseMean}{-1.83}

% ── Finding 3: Method comparison ──
\providecommand{\AppBaselineMedAPE}{21.3}
\providecommand{\AppBaselineMedAE}{0.122}
\providecommand{\AppLogitBenchRegMedAPE}{17.2}
\providecommand{\AppLogitBenchRegMedAE}{0.080}
\providecommand{\AppSVDLogitMedAPE}{16.9}
\providecommand{\AppSVDLogitMedAE}{0.076}
\providecommand{\AppBlendMedAPE}{16.4}
\providecommand{\AppBlendMedAE}{0.072}
\providecommand{\AppBlendImprovement}{23}

% ── Finding 4: Per-benchmark predictability table ──
\providecommand{\AppPerBenchmarkPredictTable}{\begin{table}[h]\small
\centering
\caption{Per-benchmark predictability (50\% holdout, 5~seeds).
  Low MedAPE = redundant; high MedAPE = unique signal.}
\label{tab:bench-predict}
\begin{tabular}{lrl}
\toprule
Benchmark & MedAPE (\%) & Status \\
\midrule
\multicolumn{3}{l}{\emph{Most redundant (easily predicted from others):}} \\
HLE & 38.5 & redundant \\
Aider Polyglot & 39.4 & redundant \\
GPQA Diamond & 40.7 & redundant \\
DeepSynth & 45.2 & redundant \\
LawBench & 46.1 & redundant \\
AIME & 49.4 & redundant \\
GSO & 50.3 & redundant \\
Widesearch & 50.5 & redundant \\
\midrule
\multicolumn{3}{l}{\emph{Most unique (hard to predict from others):}} \\
MMAU & 419.2 & unique \\
KUMO & 295.2 & unique \\
Reasoning Gym & 254.3 & unique \\
SWE-Bench Pro & 228.9 & unique \\
IneqMath & 165.5 & unique \\
HumanEvalFix & 145.9 & unique \\
SLDBench & 143.7 & unique \\
CompileBench & 139.9 & unique \\
\bottomrule
\end{tabular}
\end{table}}

% ── BenchPress top redundant / unique ──
\providecommand{\BenchPressRedundantOne}{HLE}
\providecommand{\BenchPressRedundantOneScore}{38.5}
\providecommand{\BenchPressUniqueOne}{MMAU}
\providecommand{\BenchPressUniqueOneScore}{419.2}
\providecommand{\BenchPressRedundantTwo}{Aider Polyglot}
\providecommand{\BenchPressRedundantTwoScore}{39.4}
\providecommand{\BenchPressUniqueTwo}{KUMO}
\providecommand{\BenchPressUniqueTwoScore}{295.2}
\providecommand{\BenchPressRedundantThree}{GPQA Diamond}
\providecommand{\BenchPressRedundantThreeScore}{40.7}
\providecommand{\BenchPressUniqueThree}{Reasoning Gym}
\providecommand{\BenchPressUniqueThreeScore}{254.3}

% ── Task-level SVD ──
\providecommand{\TaskTotalTasks}{6801}
\providecommand{\TaskZeroVarTasks}{908}
\providecommand{\TaskGoodTasks}{5893}
\providecommand{\TaskFullPCOne}{31.0}
\providecommand{\TaskFullPCTwo}{13.9}
\providecommand{\TaskModelPCOne}{41.8}
\providecommand{\TaskAgentPCTwoMean}{81.7}
\providecommand{\TaskBasePCTwoMean}{-81.7}

% ── Task redundancy ──
\providecommand{\TaskRedundNumBenchmarks}{52}
\providecommand{\TaskRedundMeanPCOne}{50.3}
\providecommand{\TaskRedundMedianPCsNinety}{8}
\providecommand{\TaskRedundMeanCompression}{15.9}
\providecommand{\TaskRedundMeanZeroVarPct}{12.1}
\providecommand{\TaskRedundRhoKOne}{0.807}
\providecommand{\TaskRedundRhoKThree}{0.846}
\providecommand{\TaskRedundRhoKFive}{0.864}
\providecommand{\TaskRedundRhoKTen}{0.891}
\providecommand{\TaskRedundAboveNFKOne}{1/52}
\providecommand{\TaskRedundAboveNFKThree}{5/52}
\providecommand{\TaskRedundAboveNFKFive}{8/52}
\providecommand{\TaskRedundAboveNFKTen}{20/52}
\providecommand{\TaskLowRankExOneName}{ResearchCodeBench}
\providecommand{\TaskLowRankExOneTasks}{212}
\providecommand{\TaskLowRankExOnePCOne}{78.4}
\providecommand{\TaskLowRankExOnePCs}{3}
\providecommand{\TaskLowRankExTwoName}{HumanEvalFix}
\providecommand{\TaskLowRankExTwoTasks}{164}
\providecommand{\TaskLowRankExTwoPCOne}{86.2}
\providecommand{\TaskLowRankExTwoPCs}{2}
\providecommand{\TaskLowRankExThreeName}{MMMLU}
\providecommand{\TaskLowRankExThreeTasks}{150}
\providecommand{\TaskLowRankExThreePCOne}{84.9}
\providecommand{\TaskLowRankExThreePCs}{3}
\providecommand{\TaskZeroVarExOneName}{PIXIU}
\providecommand{\TaskZeroVarExOnePct}{68}
\providecommand{\TaskZeroVarExTwoName}{StrongReject}
\providecommand{\TaskZeroVarExTwoPct}{52}
\providecommand{\TaskZeroVarExThreeName}{Spider 2}
\providecommand{\TaskZeroVarExThreePct}{44}

\providecommand{\AppGoToTaskTable}{\begin{table}[t]\small
\centering
\caption{Best single representative task per benchmark, ranked by useful representativeness score (leave-one-out correlation $\times$ task variance). Full list in \texttt{task\_representative\_tasks.csv}.}
\label{tab:go-to-tasks}
\begin{tabular}{llc}
\toprule
Benchmark & Best representative task & Score \\
\midrule
MMMLU & \texttt{mmmlu-en-us-00032} & 0.49 \\
ResearchCodeBench & \texttt{fractalgen\_unpatchify} & 0.46 \\
GAIA & \texttt{ed58682d-bc52-4baa-9eb0-4eb81e1edacc} & 0.42 \\
FinanceAgent & \texttt{financeagent-29} & 0.41 \\
LAB-Bench & \texttt{figqa-0128} & 0.41 \\
SWE-Bench Pro & \texttt{instance\_navidrome\_\_navidrome-09ae41a...} & 0.41 \\
ARC-AGI-2 & \texttt{8f215267\_0} & 0.41 \\
SWE-Bench-Multilingual & \texttt{swe-bench/swebench\_multilingual\_\_redi...} & 0.41 \\
HLE & \texttt{hle/hle\_\_671f1f4ae38f776acdad8a77} & 0.41 \\
Omni-math & \texttt{omnimath\_3346} & 0.41 \\
TerminalBench2.0 & \texttt{sqlite-db-truncate} & 0.41 \\
GSO & \texttt{gso-pandas-dev--pandas-2421931} & 0.41 \\
DeepSynth & \texttt{69} & 0.40 \\
Aider Polyglot & \texttt{polyglot\_python\_food-chain} & 0.40 \\
FeatureBench & \texttt{mwaskom\_\_seaborn.7001ebe7.test\_bar.12...} & 0.39 \\
SWE-bench-verified & \texttt{pydata\_\_xarray-6721} & 0.38 \\
GPQA Diamond & \texttt{24} & 0.38 \\
Spider 2 & \texttt{quickbooks002} & 0.38 \\
SWT Bench & \texttt{matplotlib\_\_matplotlib-24177} & 0.38 \\
USACO & \texttt{132} & 0.38 \\
SpreadsheetBench & \texttt{438-18} & 0.37 \\
PIXIU & \texttt{pixiu-mlesg-mlesg10} & 0.37 \\
SWE-smith & \texttt{oauthlib\_\_oauthlib.1fd52536.combine\_f...} & 0.37 \\
BIX-Bench & \texttt{bix-61-q2} & 0.36 \\
ReplicationBench & \texttt{disk\_ridges\_\_ridge\_slope} & 0.36 \\
\bottomrule
\end{tabular}
\end{table}}

% ── Global cross-benchmark go-to tasks ──
\providecommand{\GlobalGoToNumTasks}{5857}
\providecommand{\GlobalGoToIndependent}{395}
\providecommand{\AppGlobalGoToTaskTable}{\begin{table}[t]\small
\centering
\caption{Top 25 globally representative go-to tasks across all benchmarks, ranked by global useful representativeness (leave-one-out correlation with the full-pool aggregate $\times$ task variance). Full list in \texttt{task\_global\_representatives.csv}.}
\label{tab:global-go-to-tasks}
\begin{tabular}{llcc}
\toprule
Benchmark & Task & Score & Difficulty \\
\midrule
FinanceAgent & \texttt{financeagent-48} & 0.38 & medium \\
HLE & \texttt{hle/hle\_\_6728ba13fbd2af689fc469e5} & 0.38 & medium \\
GAIA & \texttt{ed58682d-bc52-4baa-9eb0-4eb81e1e...} & 0.38 & medium \\
SpreadsheetBench & \texttt{36277} & 0.38 & medium \\
Aider Polyglot & \texttt{polyglot\_java\_simple-linked-list} & 0.38 & medium \\
BIX-Bench & \texttt{bix-61-q2} & 0.37 & medium \\
DeepSynth & \texttt{35} & 0.37 & medium \\
PIXIU & \texttt{pixiu-mlesg-mlesg10} & 0.37 & medium \\
SpreadsheetBench & \texttt{48969} & 0.37 & medium \\
DeepSynth & \texttt{69} & 0.37 & medium \\
BIX-Bench & \texttt{bix-12-q5} & 0.37 & easy \\
FinanceAgent & \texttt{financeagent-40} & 0.37 & medium \\
SpreadsheetBench & \texttt{50811} & 0.37 & medium \\
HLE & \texttt{hle/hle\_\_66fec7825e6051260840e060} & 0.37 & medium \\
Aider Polyglot & \texttt{polyglot\_javascript\_killer-sudok...} & 0.37 & medium \\
LAB-Bench & \texttt{figqa-0024} & 0.37 & medium \\
SpreadsheetBench & \texttt{263-1} & 0.37 & easy \\
SWE-Bench Pro & \texttt{instance\_navidrome\_\_navidrome-09...} & 0.37 & easy \\
FinanceAgent & \texttt{financeagent-26} & 0.37 & medium \\
SpreadsheetBench & \texttt{14240} & 0.37 & medium \\
TerminalBench2.0 & \texttt{feal-differential-cryptanalysis} & 0.37 & medium \\
Reasoning Gym & \texttt{reasoning-gym-cognition-figlet-f...} & 0.37 & medium \\
TerminalBench2.0 & \texttt{reshard-c4-data} & 0.37 & medium \\
SpreadsheetBench & \texttt{59639} & 0.37 & easy \\
ARC-AGI-2 & \texttt{58490d8a\_0} & 0.37 & medium \\
\bottomrule
\end{tabular}
\end{table}}

\providecommand{\AppGlobalGreedySequence}{  \item Reasoning Gym / \texttt{reasoning-gym-cognition-col...} ($r{=}0.00$)
  \item Seal0 / \texttt{57} ($r{=}0.00$)
  \item Omni-math / \texttt{omnimath\_3664} ($r{=}0.00$)
  \item HLE / \texttt{hle/hle\_\_6737591afaa3cc153f...} ($r{=}0.02$)
  \item Reasoning Gym / \texttt{reasoning-gym-algorithmic-c...} ($r{=}0.06$)
  \item Aider Polyglot / \texttt{polyglot\_python\_list-ops} ($r{=}0.09$)
  \item Widesearch / \texttt{widesearch/widesearch-ws-en...} ($r{=}0.12$)
  \item SpreadsheetBench / \texttt{52305} ($r{=}0.17$)
  \item HLE / \texttt{hle/hle\_\_673511ca392d03ca70...} ($r{=}0.19$)
  \item BigCodeBench / \texttt{bigcodebench\_302} ($r{=}0.21$)}

% ── Combined go-to task set ──
\providecommand{\GoToSetNumTasks}{159}
\providecommand{\GoToSetNumBenchmarks}{53}
\providecommand{\AppGoToTaskSetTable}{\begin{table}[t]\small
\centering
\caption{Combined go-to task set: top-3 independently informative and representative tasks per benchmark (greedy $\max|r|{<}0.7$, then ranked by useful representativeness). Showing first 30 of 159 tasks across 53 benchmarks.}
\label{tab:goto-task-set}
\begin{tabular}{llccc}
\toprule
Benchmark & Task & Repr.\ score & Greedy step & Difficulty \\
\midrule
AA-LCR & \texttt{aa-lcr/aa-lcr-19} & 0.33 & 32 & easy \\
AA-LCR & \texttt{aa-lcr/aa-lcr-51} & 0.30 & 45 & easy \\
AA-LCR & \texttt{aa-lcr/aa-lcr-32} & 0.29 & 46 & easy \\
\addlinespace
Aider Polyglot & \texttt{polyglot\_javascript\_rest-api} & 0.35 & 2 & easy \\
Aider Polyglot & \texttt{polyglot\_javascript\_palindrome-p...} & 0.32 & 50 & medium \\
Aider Polyglot & \texttt{polyglot\_rust\_accumulate} & 0.32 & 47 & medium \\
\addlinespace
AIME & \texttt{aime\_i-14} & 0.21 & 8 & medium \\
AIME & \texttt{aime\_ii-11} & 0.19 & 10 & easy \\
AIME & \texttt{aime\_ii-9} & 0.19 & 11 & easy \\
\addlinespace
AlgoTune & \texttt{algotune-count-riemann-zeta-zeros} & 0.31 & 59 & medium \\
AlgoTune & \texttt{algotune-kalman-filter} & 0.23 & 66 & medium \\
AlgoTune & \texttt{algotune-ode-stiff-vanderpol} & 0.22 & 62 & medium \\
\addlinespace
ARC-AGI-2 & \texttt{7b3084d4\_0} & 0.27 & 9 & hard \\
ARC-AGI-2 & \texttt{db695cfb\_0} & 0.26 & 2 & medium \\
ARC-AGI-2 & \texttt{e3721c99\_1} & 0.23 & 10 & hard \\
\addlinespace
BFCL & \texttt{bfcl-live-simple-123-79-0} & 0.22 & 29 & medium \\
BFCL & \texttt{bfcl-live-parallel-multiple-8-7-0} & 0.19 & 36 & easy \\
BFCL & \texttt{bfcl-live-multiple-404-140-0} & 0.18 & 5 & easy \\
\addlinespace
BigCodeBench & \texttt{bigcodebench\_618} & 0.27 & 20 & medium \\
BigCodeBench & \texttt{bigcodebench\_752} & 0.25 & 2 & easy \\
BigCodeBench & \texttt{bigcodebench\_139} & 0.24 & 58 & easy \\
\addlinespace
BIX-Bench & \texttt{bix-12-q4} & 0.32 & 24 & medium \\
BIX-Bench & \texttt{bix-24-q2} & 0.29 & 27 & medium \\
BIX-Bench & \texttt{bix-41-q4} & 0.29 & 22 & easy \\
\addlinespace
CompileBench & \texttt{jq-static-musl} & 0.29 & 4 & easy \\
CompileBench & \texttt{jq-static} & 0.28 & 5 & easy \\
CompileBench & \texttt{coreutils} & 0.23 & 6 & easy \\
\addlinespace
CRUST-Bench & \texttt{crustbench-simple-lang} & 0.32 & 16 & easy \\
CRUST-Bench & \texttt{crustbench-clhash} & 0.31 & 2 & medium \\
CRUST-Bench & \texttt{crustbench-ljmm} & 0.25 & 17 & easy \\
\bottomrule
\end{tabular}
\end{table}}

\providecommand{\GoToOverlapCount}{10}
\providecommand{\GoToOverlapTotal}{54}
\providecommand{\GoToOverlapPct}{19}

% ── Unified task holdout comparison ──
\providecommand{\UnifiedHoldoutNumBenchmarks}{54}
\providecommand{\UnifiedWithinBetter}{43}
\providecommand{\UnifiedCrossBetter}{11}
\providecommand{\UnifiedWithinBetterPct}{80}
\providecommand{\UnifiedMedianWithinImp}{12.6}
\providecommand{\UnifiedMedianCrossImp}{2.6}
\providecommand{\UnifiedMedianGlobalImp}{3.3}

% @@STATS_END@@

%%%%%%%%%%%%%%%%%%%%%%%%%%%%%%%%%%%%%%%%%%%%%%%%%%%%%%%%%%%%%%%%%%%%%%%%%%%%%%%
\subsubsection{Benchmark Redundancy Analysis}
\label{app:benchpress}
%%%%%%%%%%%%%%%%%%%%%%%%%%%%%%%%%%%%%%%%%%%%%%%%%%%%%%%%%%%%%%%%%%%%%%%%%%%%%%%

We adapt the BenchPress methodology~\cite{papailiopoulos2026benchpress} to
study how much redundancy exists among our benchmarks and whether agent
scaffolding introduces structure beyond what model capability alone explains.
BenchPress was originally applied to an 83-model $\times$ 49-benchmark
matrix (34\% filled); its core analysis used a fully-observed
31-model $\times$ 10-benchmark subblock.  Our matrix is
\AppNumSystems{}~system $\times$ \AppNumBenchmarks{}~benchmark (fully observed after filtering
extreme-valued and manually excluded columns). \bo{Should we clarify the exact data object used here? In particular, it would help to state whether this is the raw benchmark matrix, the task first benchmark aggregate matrix, or the normalized benchmark relative matrix. Since later conclusions depend on the score space, I think one short 'data contract' sentence here would make the appendix much easier to audit.}

\paragraph{Methodology.}
All percentage-bounded scores $s \in [0,1]$ are transformed to logit space
before any regression or SVD:
\begin{equation}
  z = \operatorname{logit}(s) = \log\!\frac{\,\mathrm{clip}(s,\,\epsilon,\,1{-}\epsilon)\,}
       {1 - \mathrm{clip}(s,\,\epsilon,\,1{-}\epsilon)},
  \quad \epsilon = 0.005.
  \label{eq:logit}
\end{equation}\bo{Could we specify how non percentage bounded benchmarks are handled? Some benchmark columns may have scores outside $[0,1]$ or nonstandard scales. If they are excluded or normalized. Maybe we should state that explicitly before the logit transform. Otherwise readers may assume all columns are treated identically.}
The BenchPress predictor is a blend of two components:

\begin{enumerate}
\item \textbf{LogitBenchReg.} For each target benchmark~$b_j$, identify the
  $k{=}5$ other benchmarks with the highest training-set $R^2$ when predicting
  $b_j$ via univariate OLS in logit space (minimum $R^2 \ge 0.1$).  For a
  held-out system~$i$, the prediction is the $R^2$-weighted average of the
  individual univariate predictions, mapped back through the inverse logit.

\item \textbf{SVD-Logit (Soft-Impute).}  Column-wise $z$-score the
  logit-transformed matrix, initialize missing entries to~0, then iterate:
  compute rank-$r$ truncated SVD, replace \emph{only} missing entries with
  the low-rank approximation (observed entries stay pinned), and repeat until
  convergence ($\Delta_{\mathrm{rel}} < 10^{-4}$).  De-standardize and
  inverse-logit to obtain predictions.
\end{enumerate}

The final predictor is
$\hat{s} = 0.6 \times \hat{s}^{\mathrm{BenchReg}} + 0.4 \times \hat{s}^{\mathrm{SVD}}$,
following~\cite{papailiopoulos2026benchpress}.  When LogitBenchReg returns no
prediction for a cell, the method falls back to SVD-Logit alone; when both
fail, it falls back to the column mean. \bo{The current description may read like an exact reproduction of\cite{papailiopoulos2026benchpress}, but our setting has only 16 systems and includes agent scaffolds. If our implementation differs from the original, for example using univariate OLS with $R^2$-weighted averaging rather than ridge regression, is it better to call it a BenchPress-inspired adaptation?}

\paragraph{Evaluation protocol.}
We use a per-system holdout: for each of 5 random seeds, randomly hide 50\%
of each system's scores, train BenchPress on the remaining cells, and predict
the hidden ones.  The primary metric is \textbf{Median Absolute Percentage
Error} (MedAPE):
\begin{equation}
  \mathrm{MedAPE} = \operatorname{median}_{(i,j)\in\mathcal{H}}
    \frac{|\hat{s}_{i,j} - s_{i,j}|}{|s_{i,j}|} \times 100\%.
  \label{eq:medape}
\end{equation}
This matches the metric used by BenchPress, making our results directly
comparable to their reported 7\% MedAPE on the larger matrix. \bo{Both our matrices and sparsity are different (including task-first aggregation and normalized benchmark protocol). Would it be better to write it as inspired by, or something like a sanity-check style comparison? } 

\paragraph{Finding~1: Low-rank structure (SVD spectrum).}
We compute SVD in logit space on the \AppNumModelsOnly{}~system $\times$ \AppNumBenchmarks{}~benchmark
Terminus-2 submatrix (baseline scaffold only, no advanced agent).
Note that Terminus-2 is a minimal agent harness rather than a bare model evaluation,
so this submatrix is not directly equivalent to the model-only matrices studied by
BenchPress; nonetheless, it isolates the effect of varying only the base model while
holding the scaffold fixed.
The first principal component explains \AppSVDModelOnlyPCOne{}\% of the variance,
comparable to BenchPress's 71\% on their 31$\times$10 block (which used truly
scaffold-free evaluations).  The first two components
capture \AppSVDModelOnlyPCOneTwo{}\%.  This suggests that the benchmark space is dominated by a
single ``general capability'' axis even when restricting to our set of
\AppNumModelsOnly{}~frontier models under a fixed scaffold.

On the full \AppNumSystems{}-system matrix (models with and without agents), the spectrum
shifts: PC1 explains \AppSVDFullPCOne{}\% and PC2 explains \AppSVDFullPCTwo{}\% (cumulative \AppSVDFullCumTwo{}\%).
The reduction in PC1's share is explained by Finding~2.

\paragraph{Finding~2: PC2 perfectly separates advanced-scaffold from
baseline systems.}
On the \AppNumSystems{}$\times$\AppNumBenchmarks{} matrix, PC2 achieves
perfect separation between the \AppNumAgentSystems{} advanced-scaffold
systems (PC2 $> 0$; mean $= {}\AppPCTwoAgentMean{}$) and the
\AppNumBaseSystems{} Terminus-2 baseline systems (PC2 $< 0$; mean $=
{}\AppPCTwoBaseMean{}$) with zero overlap. This means the second
principal component of the combined score matrix captures the ``agent
effect''---a dimension orthogonal to general model capability that
BenchPress could not observe, since its matrix contained only baseline
systems.

\paragraph{Finding~3: BenchPress blend reduces prediction error by \AppBlendImprovement{}\%
over the column-mean baseline.}
Table~\ref{tab:benchpress-methods} compares prediction methods at 50\%
holdout.

\begin{table}[h]\small
\centering
\caption{Prediction methods on our \AppNumSystems{}$\times$\AppNumBenchmarks{} matrix (50\% holdout,
5~seeds).  MedAPE (\%) is the primary metric.}
\label{tab:benchpress-methods}
\begin{tabular}{lcc}
\toprule
Method & MedAPE (\%) & MedAE \\
\midrule
Column mean (baseline) & \AppBaselineMedAPE{} & \AppBaselineMedAE{} \\
LogitBenchReg only     & \AppLogitBenchRegMedAPE{} & \AppLogitBenchRegMedAE{} \\
SVD-Logit ($r{=}2$)    & \AppSVDLogitMedAPE{} & \AppSVDLogitMedAE{} \\
BenchPress blend ($r{=}2$) & \textbf{\AppBlendMedAPE{}} & \textbf{\AppBlendMedAE{}} \\
\bottomrule
\end{tabular}
\end{table}

The BenchPress blend achieves \AppBlendMedAPE{}\% MedAPE, a \AppBlendImprovement{}\% relative improvement
over the \AppBaselineMedAPE{}\% column-mean baseline.  This is higher than BenchPress's
7\% on its larger matrix, which we attribute to our smaller training set
(\AppNumSystems{} vs.\ 31 rows) and greater benchmark diversity (\AppNumBenchmarks{} vs.\ 10 columns).

\paragraph{Finding~4: Per-benchmark predictability.}
Table~\ref{tab:bench-predict} shows the most and least predictable benchmarks,
sorted by MedAPE at 50\% holdout.

\AppPerBenchmarkPredictTable{}

Benchmarks that resist prediction---\BenchPressUniqueOne{}
(\BenchPressUniqueOneScore{}\%), \BenchPressUniqueTwo{}
(\BenchPressUniqueTwoScore{}\%), \BenchPressUniqueThree{}
(\BenchPressUniqueThreeScore{}\%)---share a common trait: they test
capabilities orthogonal to the mainstream evaluation axis (novel agentic
research tasks, scientific experimentation, abstract reasoning).
Conversely, highly predictable benchmarks like \BenchPressRedundantOne{}
(\BenchPressRedundantOneScore{}\%) and \BenchPressRedundantTwo{}
(\BenchPressRedundantTwoScore{}\%) are fully explained by the
general-capability component.


\paragraph{Finding~5: Greedy forward selection identifies a minimal independent benchmark set.}
\label{par:greedy-selection}
To determine the smallest subset of benchmarks that captures most independent information, we apply greedy forward selection: starting from the benchmark with the lowest mean absolute correlation to all others, we iteratively add the benchmark whose maximum absolute correlation with the already-selected set is smallest.  The procedure orders all \AppNumBenchmarks{} benchmarks from most to least independent.

The first \GreedyBenchmarksBelowSeven{} benchmarks selected all have $\max|r| < 0.7$ to the selected set, meaning each adds substantial independent signal.  From step~\GreedyFirstAboveSevenStep{} onward, all remaining benchmarks have $\max|r| \ge 0.7$ with at least one already-selected benchmark---they are increasingly redundant.

The selection sequence begins with:
\begin{enumerate}
  \item CodePDE ($r{=}0.00$, Scientific Research)
  \item PIXIU ($r{=}0.00$, Professional Domains)
  \item IneqMath ($r{=}0.25$, Mathematics)
  \item ResearchCodeBench ($r{=}0.26$, Scientific Research)
  \item BFCL ($r{=}0.42$, Agents/Tools)
  \item AlgoTune ($r{=}0.45$, Software Engineering)
  \item AA-LCR ($r{=}0.56$, Software Engineering)
  \item USACO ($r{=}0.56$, Software Engineering)
  \item LAB-Bench ($r{=}0.59$, Scientific Research)
  \item SLDBench ($r{=}0.61$, Scientific Research)
\end{enumerate}
The first redundant benchmark (GSO, step~\GreedyFirstAboveSevenStep{}, $r{=}0.70$) is largely covered by the software engineering benchmarks already selected.  This result implies that, for practitioners with constrained evaluation budgets, evaluating on $\sim$\GreedyBenchmarksBelowSeven{} well-chosen benchmarks across domains can capture most of the independent signal in a \AppNumBenchmarks{}-benchmark suite.

\paragraph{Extension to the Task Level.}
While BenchPress was designed for continuous benchmark-level scores, we extend its SVD methodology to the task level (predicting individual question outcomes) to verify if the structural separation between models and agents persists at high granularity.

After filtering out \TaskZeroVarTasks{} zero-variance tasks (where all systems scored 0 or 1), we apply logit-transformed SVD to the remaining \AppNumSystems{}-system $\times$ \TaskGoodTasks{}-task matrix. Despite the high dimensionality and the binary or near-binary nature of many individual task scores, the low-rank structure holds:
\begin{itemize}
    \item \textbf{General Capability (PC1):} Explains \TaskFullPCOne{}\% of the variance across all \AppNumSystems{} systems (and \TaskModelPCOne{}\% when restricted to the \AppNumModelsOnly{} baseline systems).
    \item \textbf{Agent Axis (PC2):} Explains \TaskFullPCTwo{}\% of the variance. Even at the task level, PC2 separates advanced-scaffold systems (mean PC2 = $+\TaskAgentPCTwoMean{}$) from baseline systems (mean PC2 = $\TaskBasePCTwoMean{}$).
\end{itemize}

The contrast between the benchmark-level and task-level SVD spectra is itself informative: at the benchmark level, $\sigma_1$ captures \AppSVDFullPCOne{}\% of variance (one dominant axis), whereas at the task level $\sigma_1$ captures only \TaskFullPCOne{}\% with a much flatter decay---meaning task-level evaluation preserves multiple independent dimensions that benchmark aggregation compresses away.

\paragraph{Intra-benchmark dimensionality.}
For each of the \TaskRedundNumBenchmarks{} benchmarks with $\ge 10$ tasks, we apply PCA to the \AppNumSystems{}-system $\times$ $T$ task score matrix (after removing zero-variance tasks, which average \TaskRedundMeanZeroVarPct{}\% per benchmark).  On average, the first principal component alone explains \textbf{\TaskRedundMeanPCOne{}\%} of the variance across tasks within a benchmark, and a median of only \textbf{\TaskRedundMedianPCsNinety{} components} suffices for 90\% cumulative variance---a \textbf{\TaskRedundMeanCompression{}$\times$} compression ratio.

Some benchmarks are particularly low-rank:
\begin{itemize}
  \item \textbf{\TaskLowRankExTwoName{}} (\TaskLowRankExTwoTasks{} tasks): PC1 = \TaskLowRankExTwoPCOne{}\%, only \TaskLowRankExTwoPCs{}~components for 90\%.
  \item \textbf{\TaskLowRankExThreeName{}} (\TaskLowRankExThreeTasks{} tasks): PC1 = \TaskLowRankExThreePCOne{}\%, \TaskLowRankExThreePCs{}~components for 90\%.
  \item \textbf{\TaskLowRankExOneName{}} (\TaskLowRankExOneTasks{} tasks): PC1 = \TaskLowRankExOnePCOne{}\%, \TaskLowRankExOnePCs{}~components for 90\%.
\end{itemize}
These benchmarks contain hundreds of items that essentially measure the same underlying capability axis.  We can also verify this via ranking fidelity: selecting the $k$ tasks with the highest Spearman correlation to the full benchmark mean score, the resulting $k$-task subset preserves the system ranking with high fidelity (Table~\ref{tab:task-redundancy}).

\begin{table}[t]
  \centering
  \caption{System ranking fidelity ($\rho$) from a small task subset vs.\ the full benchmark.  Values are mean Spearman rank correlation across all benchmarks with $\ge 10$ tasks.}
  \label{tab:task-redundancy}
  \begin{tabular}{cccc}
    \toprule
    $k$ tasks & Mean $\rho$ & Benchmarks $\rho > 0.95$ & Typical reduction \\
    \midrule
    1  & \TaskRedundRhoKOne{}  & \TaskRedundAboveNFKOne{}  & $>$100$\times$ \\
    3  & \TaskRedundRhoKThree{} & \TaskRedundAboveNFKThree{} & 30--60$\times$ \\
    5  & \TaskRedundRhoKFive{}  & \TaskRedundAboveNFKFive{}  & 20--40$\times$ \\
    10 & \TaskRedundRhoKTen{}   & \TaskRedundAboveNFKTen{}   & 10--20$\times$ \\
    \bottomrule
  \end{tabular}
\end{table}

With just 10 well-chosen tasks, \TaskRedundAboveNFKTen{} benchmarks achieve $\rho > 0.95$, and the overall mean $\rho$ across all benchmarks reaches \TaskRedundRhoKTen{} at $k{=}10$---indicating that a small task subset can approximate the full-suite ranking with high fidelity for most benchmarks.

\paragraph{Within-benchmark task predictability via holdout.}
\label{par:task-holdout}
The PCA and ranking analyses above quantify redundancy from spectral and rank-correlation perspectives.  We now apply the BenchPress prediction framework directly: can we predict hidden task scores from peer tasks?

While applying the full BenchPress pipeline (LogitBenchReg from $\sim$6000 tasks) to the task level fails due to dimensionality---$\TaskGoodTasks{}$ tasks with only $\AppNumSystems{}{-}1$ training systems guarantees spurious $R^2{=}1$---we circumvent this by restricting prediction to tasks \emph{within the same benchmark}.

For each benchmark with $\ge 5$ reliable tasks, we apply a BenchPress-style holdout protocol within the benchmark's $(16 \times T)$ sub-matrix:
\begin{enumerate}
  \item For each fold (3 folds), randomly hide 50\% of each system's task scores.
  \item Predict hidden cells using three methods:
    \textbf{(a)} task-mean baseline (column mean of visible systems),
    \textbf{(b)} top-$k$ peer weighted (correlation-weighted average of the $k$ most correlated peer tasks), and
    \textbf{(c)} within-benchmark rank-2 SVD (soft-impute).
  \item Blend peer + SVD predictions ($0.6 \times \text{peer} + 0.4 \times \text{SVD}$).
  \item Report per-task Median Absolute Error (MedAE).
\end{enumerate}

Across 53 benchmarks, the blend improves over the task-mean baseline in 48 (median improvement 28.4\%).  Benchmarks with the most internally predictable tasks---HumanEvalFix (71\% improvement), KUMO (55\%), LiveCodeBench (55\%)---contain highly redundant items that test essentially the same capability axis.  The few benchmarks where the blend underperforms the baseline (e.g., LawBench, SWE-Lancer) tend to have heterogeneous task structures that resist peer-based prediction.

\paragraph{Per-benchmark representative tasks and greedy selection.}
\label{par:per-bench-goto}
The holdout analysis identifies \emph{which} tasks are predictable; we now identify which tasks are most \emph{representative}.  Table~\ref{tab:go-to-tasks} lists the single most representative task per benchmark---the task whose leave-one-out correlation with the remaining-task aggregate is highest, weighted by task variance (the ``useful representativeness score'').  For practitioners who can run only one task per benchmark, these items maximize ranking fidelity to the full suite; with $k{=}3$ per benchmark the median ranking fidelity reaches $\rho \approx 0.89$.

\AppGoToTaskTable{}

Mirroring the benchmark-level greedy forward selection (Finding~5), we also apply greedy selection within each benchmark: iteratively select the task whose maximum absolute correlation with already-selected tasks is smallest.  The fraction of tasks with $\max|r| < 0.7$ varies widely: benchmarks like BigCodeBench (58\%) and AlgoTune (56\%) contain many independently-informative items, while highly redundant benchmarks like HumanEvalFix and MMMLU reach $r > 0.7$ within the first few selections.

\paragraph{Unified within- vs.\ cross-benchmark task prediction.}
\label{par:unified-holdout}
The analyses above use only within-benchmark peers.  A natural question is whether tasks from \emph{other} benchmarks can predict a given task better than same-benchmark peers.  To answer this fairly, we run a unified holdout experiment: for each benchmark, we apply the \emph{same} holdout mask and the \emph{same} $k{=}5$ peer count across three peer scopes:
\begin{itemize}
  \item \textbf{Within-benchmark peers:} the $k$ most-correlated tasks from the same benchmark.
  \item \textbf{Cross-benchmark peers:} the $k$ most-correlated tasks from \emph{other} benchmarks only.
  \item \textbf{Global peers:} the $k$ most-correlated tasks from the entire pool (any benchmark).
\end{itemize}
Each scope produces a peer prediction and a blend with rank-2 SVD ($0.6 \times \text{peer} + 0.4 \times \text{SVD}$).  We report MedAE improvement over the task-mean baseline.

Across \UnifiedHoldoutNumBenchmarks{} benchmarks, within-benchmark blends outperform global blends in \UnifiedWithinBetter{}/\UnifiedHoldoutNumBenchmarks{} (\UnifiedWithinBetterPct{}\%) cases.  The median improvement over baseline is \UnifiedMedianWithinImp{}\% (within), \UnifiedMedianCrossImp{}\% (cross-benchmark), and \UnifiedMedianGlobalImp{}\% (global).  For the majority of benchmarks, same-benchmark peers provide more accurate predictions---confirming that tasks within a benchmark genuinely share structure beyond what cross-benchmark signal captures.  However, for a substantial minority (\UnifiedCrossBetter{}/\UnifiedHoldoutNumBenchmarks{}), cross-benchmark or global peers predict \emph{better}, suggesting that some benchmarks share latent capability axes that cut across evaluation boundaries.

\paragraph{Cross-benchmark global go-to tasks.}
\label{par:global-goto}
The per-benchmark representative analysis (Table~\ref{tab:go-to-tasks}) identifies the best task \emph{within} each benchmark, but practitioners may want a single compact task set that covers as much independent signal as possible \emph{across all benchmarks}.  We flatten all \GlobalGoToNumTasks{} reliable bounded tasks into a single pool and perform two complementary analyses:

\begin{enumerate}
  \item \textbf{Global representativeness.}  For each task, we compute the leave-one-out Pearson correlation between its normalized score profile across all 16 systems and the aggregate profile of all remaining tasks in the pool.  Multiplying by the task's observed standard deviation yields a ``global useful representativeness'' score---high values indicate tasks that both track the overall system ranking and have enough variance to discriminate systems.

  \item \textbf{Greedy forward selection across all tasks.}  Starting from the task with the lowest mean absolute Pearson correlation to all others, we iteratively add the task whose maximum absolute correlation with the already-selected set is smallest.  This orders tasks from most to least independent.
\end{enumerate}

Table~\ref{tab:global-go-to-tasks} lists the top~25 tasks by global useful representativeness.  The greedy selection sequence begins with:
\begin{enumerate}
\AppGlobalGreedySequence{}
\end{enumerate}

\AppGlobalGoToTaskTable{}

\paragraph{Per-benchmark vs.\ global representative overlap.}
A key question is whether the per-benchmark ``best representative task'' and the globally most representative tasks are the same set.  Comparing the per-benchmark top-1 tasks (one per benchmark) with the top-$N$ globally representative tasks (where $N$ equals the number of benchmarks), we find only \GoToOverlapCount{}/\GoToOverlapTotal{} (\GoToOverlapPct{}\%) overlap.
This is because per-benchmark representativeness selects tasks that best track \emph{within-benchmark} variation, while global representativeness selects tasks that best track the \emph{cross-benchmark consensus}.  Benchmarks with idiosyncratic score profiles (e.g., those identified as ``unique'' in the BenchPress analysis) nominate internally representative tasks that look unrepresentative in the global pool.

\paragraph{Combined go-to task set.}
For practitioners seeking a compact evaluation set, we construct a combined go-to task set that balances independence and representativeness.  Within each benchmark, we apply greedy forward selection to identify independently informative tasks ($\max|r| < 0.7$), then select the top 3 by useful representativeness score.  This yields \GoToSetNumTasks{} tasks across \GoToSetNumBenchmarks{} benchmarks---a substantial reduction from the full \GlobalGoToNumTasks{}-task pool, while preserving both within-benchmark coverage and cross-benchmark diversity.

\AppGoToTaskSetTable{}

\paragraph{Implications.}
The two-level redundancy---across benchmarks (Findings~4--5) and within benchmarks (this section)---implies that a carefully selected subset of benchmarks, each evaluated on a small number of representative tasks, can approximate the full evaluation matrix at a fraction of the computational cost.  For practitioners with constrained evaluation budgets, the combined go-to task set provides a principled starting point: \GoToSetNumTasks{} tasks that balance within-benchmark representativeness with cross-benchmark independence.

%%%%%%%%%%%%%%%%%%%%%%%%%%%%%%%%%%%%%%%%%%%%%%%%%%%%%%%%%%%%%%%%%%%%%%%%%%%%%%%

\subsubsection{Model vs.\ Agent Effect: Methodology Details}
\label{app:model-vs-agent}

This appendix formalizes the within-family and adjusted-effect analyses
summarized in Section~\ref{sec:agent-harness}.  Let $s(m, a, b)$ denote the
benchmark score for model~$m$, agent scaffold~$a$, and benchmark~$b$.

%%%%%%%%%%%%%%%%%%%%%%%%%%%%%%%%%%%%%%%%%%%%%%%%%%%%%%%%%%%%%%%%%%%%%%%%%%%%%%%
\paragraph{Within-Family Comparison}
\label{app:within-family}

For each model provider family $f$ (e.g., Anthropic, OpenAI, Google), let
$M_f$ be the set of models belonging to $f$ and $A_f$ the set of non-baseline
agents evaluated on that family.  The baseline agent is Terminus~2.

\subparagraph{Model effect (fixing agent).}
For each benchmark~$b$, we fix the agent to Terminus~2 and measure the score
range across models within the family:
\begin{equation}
  \Delta_{\mathrm{model}}^{(f,b)}
  \;=\;
  \max_{m \in M_f}\, s(m,\,\text{terminus-2},\,b)
  \;-\;
  \min_{m \in M_f}\, s(m,\,\text{terminus-2},\,b).
  \label{eq:model-effect}
\end{equation}
This requires at least two models in~$f$ with observed Terminus~2 scores on~$b$.

\subparagraph{Agent effect (fixing model).}
For each model $m \in M_f$ evaluated under both Terminus~2 and another agent
$a \in A_f$ on benchmark~$b$, we compute the absolute score change:
\begin{equation}
  \delta_{\mathrm{agent}}^{(m,a,b)}
  \;=\;
  \bigl|\, s(m,\,a,\,b) - s(m,\,\text{terminus-2},\,b) \,\bigr|.
  \label{eq:agent-delta}
\end{equation}
The family-level agent effect on~$b$ averages over all available model--agent pairs:
\begin{equation}
  \Delta_{\mathrm{agent}}^{(f,b)}
  \;=\;
  \frac{1}{|\mathcal{P}_{f,b}|}
  \sum_{(m,a)\in \mathcal{P}_{f,b}}
  \delta_{\mathrm{agent}}^{(m,a,b)},
  \label{eq:agent-effect}
\end{equation}
where $\mathcal{P}_{f,b} = \{(m,a) : m \in M_f,\; a \in A_f,\; \text{both }
s(m,a,b) \text{ and } s(m,\text{terminus-2},b) \text{ are observed}\}$.

\subparagraph{Dominance classification.}
For a given $(f, b)$ pair, the benchmark is classified as
\emph{model-dominant} if $\Delta_{\mathrm{model}}^{(f,b)} >
\Delta_{\mathrm{agent}}^{(f,b)}$, and \emph{agent-dominant} otherwise.

\subparagraph{Family-level summary.}
We aggregate across benchmarks within each family using the mean:
\begin{align}
  \overline{\Delta}_{\mathrm{model}}^{(f)}
  &= \frac{1}{|B_f|}\sum_{b \in B_f} \Delta_{\mathrm{model}}^{(f,b)}, &
  \overline{\Delta}_{\mathrm{agent}}^{(f)}
  &= \frac{1}{|B_f|}\sum_{b \in B_f} \Delta_{\mathrm{agent}}^{(f,b)},
  \label{eq:family-summary}
\end{align}
where $B_f$ is the set of benchmarks with both model-effect and agent-effect
estimates for family~$f$.  The model-to-agent ratio is
$\overline{\Delta}_{\mathrm{model}}^{(f)} / \overline{\Delta}_{\mathrm{agent}}^{(f)}$,
and the model-dominant percentage is the fraction of benchmarks in~$B_f$
classified as model-dominant.

Table~\ref{tab:within-family-summary} reports these statistics for the three
multi-model families.

\begin{table}[t]
  \centering
  \caption{Within-family model vs.\ agent effect comparison.  The model effect
    is the score range across models within the family on Terminus~2; the agent
    effect is the mean absolute score change from switching the agent while
    holding the model fixed.  Ratio and model-dominant percentage are computed
    per benchmark, then averaged.}
  \label{tab:within-family-summary}
  \begin{tabular}{lcccccc}
    \toprule
    Family
      & Models
      & Agents
      & $\overline{\Delta}_{\mathrm{model}}$
      & $\overline{\Delta}_{\mathrm{agent}}$
      & Ratio
      & Model-dom.\ \% \\
    \midrule
    Anthropic
      & 3
      & claude-code
      & \AnthropicModelEffectMean{}
      & \AnthropicAgentEffectMean{}
      & \AnthropicModelAgentRatio{}$\times$
      & \AnthropicModelWinPct\% \\
    OpenAI
      & 3
      & codex
      & \OpenAIModelEffectMean{}
      & \OpenAIAgentEffectMean{}
      & \OpenAIModelAgentRatio{}$\times$
      & \OpenAIModelWinPct\% \\
    Google
      & 2
      & gemini-cli
      & \GoogleModelEffectMean{}
      & \GoogleAgentEffectMean{}
      & \GoogleModelAgentRatio{}$\times$
      & \GoogleModelWinPct\% \\
    \bottomrule
  \end{tabular}
\end{table}

Table~\ref{tab:within-family-example} illustrates the per-benchmark comparison
for one representative family.

\begin{table}[t]
  \centering
  \caption{Example within-family comparison for the Anthropic family
    (claude-opus-4-6, claude-sonnet-4-6, claude-haiku-4-5).
    \emph{Model effect}: score range across the three models on Terminus~2.
    \emph{Agent effect}: mean $|\Delta|$ from switching to Claude~Code while
    holding each model fixed.  Only a representative subset of benchmarks is
    shown.}
  \label{tab:within-family-example}
  \small
  \begin{tabular}{lcccc}
    \toprule
    & \multicolumn{2}{c}{Model effect (Terminus-2)}
    & \multicolumn{2}{c}{Agent effect (fix model)} \\
    \cmidrule(lr){2-3}\cmidrule(lr){4-5}
    Benchmark & Score range & Rank & Mean $|\Delta|$ & Rank \\
    \midrule
    \multicolumn{5}{c}{\emph{Auto-populated from pipeline output ---
      see \texttt{benchmark\_within\_family\_model\_vs\_agent.csv}}} \\
    \bottomrule
  \end{tabular}
\end{table}


